\documentclass[10pt]{beamer}

% --- 删除或注释掉以下两行 ---
% \usepackage[T2A]{fontenc} 
% \usepackage[utf8]{inputenc}

% --- 替换为以下配置 ---
\usepackage[russian]{babel}
\usepackage{fontspec}

% 指定支持俄语的系统字体 (确保你电脑上有这些字体)
\setmainfont{Times New Roman}
\setsansfont{Arial}
\setmonofont{Courier New}

\usetheme{Madrid}

\usecolortheme{default}

% --- 文档信息 ---
\title{Моя первая презентация в LaTeX}
\subtitle{Изучение структуры Beamer}
\author{Ли Хан}
\institute{Российский Университет Дружбы Народов}
\date{\today}

\begin{document}
	
	% 1. 标题页 (Title Frame)
	\begin{frame}
		\titlepage
	\end{frame}
	
	% 2. 目录页 (Table of Contents)
	\begin{frame}{Содержание}
		\tableofcontents
	\end{frame}
	
	\section{Введение}
	% 3. 普通内容页
	\begin{frame}{Цель работы}
		В этом разделе мы изучаем структуру презентации:
		\begin{itemize}
			\item Использование класса \texttt{beamer}.
			\item Создание отдельных кадров через \texttt{frame}.
			\item Организация логических разделов через \texttt{section}.
		\end{itemize}
	\end{frame}
	
	\section{Основная часть}
	% 4. 带有强调块的页面
	\begin{frame}{Важные определения}
		\begin{block}{Определение структуры}
			Презентация состоит из последовательности окружений \texttt{frame}.
		\end{block}
		
		\begin{alertblock}{Внимание!}
			Каждая страница PDF соответствует одному состоянию кадра.
		\end{alertblock}
	\end{frame}
	
\end{document}